% 原文 pp.216–222（印刷 pp.1334–1340）
% 承接第 15 讲第 1 部分：本文件自 §15.2 末段（p.216）起，经 §15.3–§15.8，至 p.222 的「参考文献与交叉引用」止；
% 因属同一讲的后续分段，故不写 \section，直接从正文与 \subsection 开始。
% 本 7 页内没有任何 Figure / 示意图，故无插图。

% ===== 原文 p.216 =====
% （本页首句承接 p.215 末段：关于此类 Calabi--Yau 流形的椭圆亏格与 Landau--Ginzburg 侧极小模型特征标的比较。）

椭圆亏格（elliptic genus）应当与 $t$ 无关，因为它是拓扑不变量；由此可以得到一个公式，把这一类 Calabi--Yau 流形的椭圆亏格，与 Landau--Ginzburg 一侧可由极小模型（minimal model）的特征标（character）显式算出的椭圆亏格联系起来。

尽管椭圆亏格与 $t$ 无关，我们讨论过的 $\mathcal{N}=2$ 理论中的另一个量——$Q_t$-上同调（$Q_t$-cohomology）——却随 $t$ 全纯地变化。这个问题中的瞬子（instanton），也就是 $\sigma$ 模型中的全纯曲线（holomorphic curve），会给路径积分贡献一个因子
\[
  \exp\Bigl(-i\theta\int\frac{F}{2\pi}-r\int d^2x\,D\Bigr),
\]
其中 $D=\sum_j\alpha_j|\Phi_j|^2-r$。$\mathcal{N}=2$ 模型的瞬子方程——即场在某个超对称下保持不变的条件——给出 $\int d^2x\,D=\int F/2\pi$，于是这个因子就是 $\exp\bigl(-2\pi it\int_\Sigma\Phi^*(c_1(\mathcal{L}))\bigr)$。我们将在本讲的后半部分探讨这些瞬子求和。

% [Tutor Note: 这里两处 $r$ 项相差一个 $2\pi$，是极易踩空之处。由 $\int d^2x\,D=\int F/2\pi=\int_\Sigma\Phi^*(c_1(\mathcal{L}))=n$，本页显示公式给出 $\exp(-irn-i\theta n)$；而 p.220 把复参数取为 $t=-ir+\theta/2\pi$，于是 $\exp(-2\pi itn)=\exp(-2\pi irn-i\theta n)$，二者在 $r$ 项上差因子 $2\pi$。要使两式一致，需把 FI 参数以 $1/2\pi$ 为单位（即 $t=\theta/2\pi-ir/2\pi$），或等价地认为 $D$ 项积分自带 $1/2\pi$ 归一化。核对第 13 讲线性 $\sigma$ 模型公式时务必先固定这一归一化，否则 p.221 的幂等元位置 $\sigma=\mu e^{2\pi i(-t+k)/n}$ 会整体错位；这与 p.221 脚注 5 所讨论的「物理归一化 vs.\ 拓扑归一化」属同一类问题。]

\subsection{量子上同调的引入}
% 原文 §15.3：Introduction to quantum cohomology

现在回到二维的扭曲（twisted）$\mathcal{N}=2$ 超对称 $\sigma$ 模型，其靶空间（target）是 Kähler 流形 $X$。回忆一下：依所保留的超对称荷 $Q\in\{\overline{Q}_+ + Q_-,Q_+ + Q_-,Q_+ + \overline{Q}_-,\overline{Q}_+ + \overline{Q}_-\}$ 的不同，共有四种扭曲方式。

% [Tutor Note: 扫描件（p.1334）中集合里几个 $Q$ 的上划线极细、几乎无法分辨，字面上会出现同一元素重复三次的无意义结果。按上下文——四种扭曲对应 $(a,a),(c,a),(a,c),(c,c)$ 四个手征环，且 p.217 明确说「今天只考虑使 $\overline{Q}_+ + Q_-$ 成为整体的扭曲」（情形 A，即 $(a,a)$ 环）——这四项应是 $\overline{Q}_\pm$ 与 $Q_\pm$ 的两两配对，译文按此恢复。若以扫描件字面为准逐字转写，须自行把缺失的上划线补回，否则后面谈「$(a,a)$ 环由 $Q=\overline{Q}_+ + Q_-$ 定义」时会自相矛盾。]

无论选择哪一种，$Q$ 的上同调作为向量空间都可与 $X$ 的上同调等同。但其中还有更多结构：若把 Riemann 面固定为某个二维流形 $\Sigma$，则额外的结构是 $Q$ 的上同调上的一个环结构。所得的环分别称为 $(a,a)$、$(c,a)$、$(a,c)$ 与 $(c,c)$ 手征环（chiral ring）。

在其中一个选择下，手征环的经典版本恰好就是上同调上通常的环结构；但在量子力学层面，对这个经典答案还有修正，即瞬子修正。正因如此，$\sigma$ 模型的手征环被称为 Kähler 流形 $X$ 的量子上同调（quantum cohomology）。在本讲的后半部分，我们将一般地讨论量子上同调，然后在 $X=\mathbb{CP}^{n-1}$ 以及 $X$ 为 Fano 超曲面的情形做一些显式计算。这些计算既从数学上做——即数具有某些性质的全纯有理曲线；也从物理的观点做：利用重整化群流（在第 13 讲与本讲前半部分讨论过），把这些 $\sigma$ 模型流向带有额外真空的 Landau--Ginzburg 理论。

\subsection{零能态的空间}
% 原文 §15.4：The space of zero-energy states

我们想计算其中一些手征环。但在直接处理算符之前，我们先在 Hamilton 框架下计算零能态（zero-energy states）。之所以要做这个计算，是因为如第 14 讲 §14.3.3 所示，存在一个映射
\[
  \psi:\ \text{手征环}\longrightarrow\{\text{零能态}\},
\]
它由在半球面上某点插入一个局域算符给出。在有利的情形下，这个映射是同构。

% ===== 原文 p.217 =====

在 Hamilton 框架下，场空间底下的玻色子空间是
\[
  \mathcal{W}=\operatorname{Maps}(S^1,X).
\]
能量最低的位形是常值映射；这些位形的空间就是 $X$ 的一份拷贝。由于旋量丛（spin bundle）的两份拷贝中各有费米子，量子化之后所得的 Hilbert 空间是 $X$ 上旋量丛两份拷贝的张量积，亦即 $X$ 上的微分形式空间。Hamiltonian 就是 Laplacian，于是它的基态空间是 $X$ 上的 $L^2$ 调和形式空间，后者与 $H^*_{L^2}(X)$ 等同。当然，还要考虑 $\mathcal{W}$ 中这份 $X$ 的法方向；但超对称意味着法方向上的二次型非退化，于是这些方向给出谐振子的张量积，并具有唯一的基态。因此在最后我们发现：整个理论的基态与 $X$ 上的调和形式等同。我们将看到，手征环的计算（在加法的意义上）给出同样的答案，于是在今天所考虑的情形里，映射 $\psi$ 确实是同构。一般地，只要经典理论具有「能量随场趋于无穷而增长」这一性质，上述结论就成立。\footnote{对于靶空间 $X$ 紧致的 $\sigma$ 模型，一个映射 $\Phi:S^1\to X$「趋于无穷」的唯一含义只是它可以沿 $S^1$ 变化得极为剧烈；这当然会使能量增长。（原文 p.217 页脚注）}

% [Tutor Note: 脚注原文作 ``the only sense in which a map $\Phi:S^1\to X$ can `go to infinity' is that \emph{they} can become rapidly varying along $S^1$''，主语 a map 与代词 they 不一致，应为 it，属排印疏漏。物理上要紧的是后半句：靶空间紧致时不存在「场趋于无穷」的通常含义，唯一可能是 $\Phi$ 沿 $S^1$ 剧烈振荡，此时能量 $\int_{S^1}|d\Phi|^2$ 确实发散。这正是保证 $\psi$ 为同构、并使超对称量子力学中 Witten 指标论证成立的关键（参见 E. Witten, J.\ Diff.\ Geom.\ 17 (1982) 661 关于超对称量子力学与 Morse 理论的工作）；若靶空间非紧，这一论证失效，$\psi$ 一般只是单射。]

\subsection{手征环的一般性质}
% 原文 §15.5：Generalities on the chiral ring

为定义手征环，我们需要一个平方为零的整体超对称。如第 14 讲 §14.3 所见，选取 $\mathcal{N}=2$ 超对称理论的某种扭曲形式，就会选出这样一个整体超对称。今天我们将只考虑使 $\overline{Q}_+ + Q_-$ 成为整体超对称的扭曲；这就是第 14 讲中所说的「情形 A」（Case A）。这种扭曲对任何 Kähler 流形都可施行——无需假定第一 Chern 类为零。我们将考虑这样的局域算符：它们是基本玻色场 $\Phi:\Sigma\to X$（$\Sigma$ 是我们的 Riemann 面）以及由超对称与 $\Phi$ 相联系的费米子的函数。（若在 $X$ 上取局部坐标，则我们用 $S^1$ 上的一族函数 $\phi_j$ 来描述 $\Phi$。）回忆扭曲的效应之一，是使费米子成为整自旋而非半自旋丛的截面；也就是说，它们不再是旋量，而是 $\Sigma$ 上的微分形式（取值于 $X$ 切丛的拉回）。在描述算符的 $Q$-上同调时，只需考虑零自旋的费米子。理由在于：所有 $Q$-上同调类的算符都有仅用这些费米子构造的代表元。更确切地说，$Q$-上同调类具有只依赖玻色场与零自旋费米子、而不依赖其导数的代表元。\footnote{这是如下事实的算符版本：在 Hamiltonian 描述中，零能态由靶空间上的微分形式构造，其中所有振子都处于基态；含场导数的算符则对应于某些振子被激发的态。（原文 p.217 页脚注）} 于是，我们的局域算符将具有形式
\[
  \mathcal{O}=\mathcal{O}(\phi,\psi_L,\psi_R),
\]
其中，在未扭曲的版本中，$\psi_L$ 是 $K^{-1/2}\otimes\phi^*T^{0,1}X$ 的截面，而 $\psi_R$ 是 $K^{1/2}\otimes T^{1,0}X$ 的截面。扭曲之后，这些费米子场位于 $\Phi^*T^{0,1}X$ 与 $\Phi^*T^{1,0}X$ 中，因为扭曲把 $\Sigma$ 上的线丛抵消掉了。特别地，

% [Tutor Note: 此处原文的 $\psi_R$ 写作 $K^{1/2}\otimes T^{1,0}X$，漏掉了拉回 $\phi^*$；紧接的「扭曲后位于 $\Phi^*T^{0,1}X$ 与 $\Phi^*T^{1,0}X$」一句可证应为 $K^{1/2}\otimes\phi^*T^{1,0}X$。另注意 $\psi_L$ 与 $\psi_R$ 携带的 $K$ 幂次相反（$-1/2$ 与 $+1/2$），这并非笔误：正因二者相反，同一个扭曲才能把它们同时抵消成 $\Phi^*T^{0,1}X$、$\Phi^*T^{1,0}X$ 的截面；若写成相同的幂次，扭曲后必残留一个 $\Sigma$ 上的线丛，$Q$ 就不再是 $\Sigma$ 上的标量超荷，$Q$-上同调也不再等于 $X$ 的 de Rham 上同调。再者，括号里说「用 $S^1$ 上的函数描述 $\Phi$」疑为 $\Sigma$ 之误：$\Phi$ 是 $\Sigma\to X$ 的映射，p.218 的 $\mathcal{O}_\lambda(x)=f_\lambda(\phi(x),\psi(x))$ 中 $x$ 属于 $\Sigma$；$S^1$ 的说法只适用于本页的 Hamiltonian 图像。]

% ===== 原文 p.218 =====

就可以组成和 $\psi=\psi_L+\psi_R$，并考虑算符 $\mathcal{O}=\mathcal{O}(\phi,\psi)$，其中 $\psi$ 是 $\Sigma$ 上的零形式，取值于 $TX=T^{1,0}X\oplus T^{0,1}X$ 的拉回。

\subsubsection{局域算符}
% 原文 §15.5.1：Local Operators

$X$ 上的每个微分形式 $\lambda$ 都对应一个函数 $f_\lambda(\phi,\psi)$（$f_\lambda$ 是 $\Pi TX$ 上的函数）。于是对每个 $\lambda$ 都有一个算符 $\mathcal{O}_\lambda(x)=f_\lambda(\phi(x),\psi(x))$。正如上周那一讲所提到的，在算符与形式的等同之下，$Q$ 等同于通常的外微分：
\[
  \{Q,\mathcal{O}_\lambda\}=\mathcal{O}_{d\lambda}.
\]
这就是为什么在加法的意义上，$Q$-上同调与 $X$ 的 de Rham 上同调一致。

% [Tutor Note: 这里的等同 $Q\leftrightarrow d$ 是整套「物理 $\leftrightarrow$ 数学」词典的核心，符号上易错者有两点。(i) $Q$ 与 $\mathcal{O}_\lambda$ 都是费米型算符（$\lambda$ 为奇次时 $\mathcal{O}_\lambda$ 亦然），故用反对易子 $\{Q,\mathcal{O}_\lambda\}=\mathcal{O}_{d\lambda}$；若误写成对易子，或忽略了 $d\lambda$ 前的符号随 $\deg\lambda$ 的变化，整体会差一个 $(-1)^{|\lambda|}$。(ii) $f_\lambda$ 是 $\Pi TX$（切丛的奇（parity）版本，纤维坐标 $\psi$ 反对易）上的函数，即「$\psi$ 的多项式 $=$ 微分形式」；而 $\lambda$ 的 de Rham 类与 Poincaré 对偶闭链 $H_i$ 的配对约定，直接决定 p.218–219 瞬子计数中 $\pm1$ 的符号（等价于 determinant line bundle 的定向选取），不同文献的符号可能整体相反。]

当算符—态对应是同构时——正如在这些 $\mathcal{N}=2$ 超对称 $\sigma$ 模型的情形——我们可以通过计算关联函数来计算手征环的环结构。这里的目标是用更经典的拓扑与几何语言来计算关联函数
\[
  \langle \mathcal{O}_{\lambda_1}(x_1)\cdots\mathcal{O}_{\lambda_s}(x_s)\rangle,
\]
其中 $\lambda_i$ 是闭形式。我们已经指出，若把 $\lambda_i$ 换成上同调等价的形式，关联函数不变。此外，要计算这个关联函数，也就是计算路径积分
\[
  \int \mathcal{D}\phi\,\mathcal{D}\psi\,e^{-\mathcal{L}}\prod_i\mathcal{O}_{\lambda_i}(x_i)
\]
（积分遍历所有场），只需把它化为在 $Q$-不变场子空间上的积分即可。我们已经辨认出 $Q$-不动点子空间的玻色部分：它就是全纯映射 $\Sigma\to X$ 的空间。对每个 $i$ 取一个与 $[\lambda_i]$ Poincaré 对偶的闭链 $H_i$；于是可以选取 $\lambda_i$ 的一个代表元，使其支集落在 $H_i$ 的任意小邻域内。

现在我们要通过把路径积分定域化（localization）到 $Q$-不变的经典位形上、即定域化到 Riemann 面到 $X$ 的全纯映射上，来计算关联函数。给定我们所选的算符代表元，只有那些在点 $x_i$ 处穿过 $H_i$ 的全纯映射才有贡献。在「最好」的情形，即过这些点的全纯映射的模空间由有限个非退化点组成时，可以证明每个这样的映射贡献 $\pm1$。事实上，超对称保证 Bose 行列式与 Fermi 行列式相消，只剩一个符号。于是，对所有全纯映射求和所得的结果，等于经典作用量的指数因子与来自行列式的符号之积。对全纯映射而言，经典作用量的指数因子恰好是 $\exp\bigl(-2\pi it\int_\Sigma\Phi^*(c_1(\mathcal{L}))\bigr)$。

% [Tutor Note: 「每个非退化映射贡献 $\pm1$」是高度压缩的结论，其依据是 Atiyah--Bott 定域化（等价地，超对称量子力学中的 Witten 不动点/指标公式）：路径积分在 $Q$-不变位形附近展开时，Bose 与 Fermi 的非零模成对出现、行列式之比为 $\pm1$，只剩符号。要定出符号，须为每个分支的 determinant line bundle 指定定向——本段「up to sign」正是回避了这一约定。还需注意：若模空间含非平凡自同构（例如恰有额外对称的映射），「有限个非退化点」的前提不成立，此时必须改用等变（equivariant）版本或虚拟基本类计数，这正是后来 Gromov--Witten 不变量的出发点（参见 Kontsevich 关于稳定映射与虚模空间的工作）。]

从数学上说，要使「在 $x_i$ 处穿过 $H_i$ 的全纯映射」这一问题给出零维模空间，必须满足一个维数计数。更一般地陈述：若把算符 $\mathcal{O}_{\lambda_i}(x_i)$ 解释为在 $\Sigma\to X$ 的所有全纯映射的模空间上产生微分形式，则 $H_i$ 须满足一个维数条件，该微分形式才是最高次形式；否则其积分必然为零。从物理观点看，同样的维数限制来自 $R$ 对称性的反常：路径积分中的积分测度在该对称性下有反常，因此必须选用这样一组算符，使其乘积

% ===== 原文 p.219 =====

恰具有可以抵消的反常，路径积分才有非零的可能。毫不奇怪，这两个条件是同一个条件：即 $X$ 中诸 $H_i$ 的余维（codimension）之和，必须等于我们所考虑的那个全纯映射模空间分支的虚拟维数（virtual dimension）。

若只想看到经典的上同调环而非量子上同调，则应取 $\Sigma$ 亏格为零，并且只考虑映入 $X$ 的常值映射。当然，常值映射恰好填满 $\operatorname{Maps}(\Sigma,X)$ 的某一分支中的全部全纯映射，即由零伦（null-homotopic）映射构成的那个分支。对这个同伦类中的映射，维数条件是
\[
  \sum_{i=1}^{s}\deg\lambda_i=\dim X,
\]
而经典答案就是
\[
  \int_X\lambda_1\wedge\cdots\wedge\lambda_s.
\]
也就是说，经典手征环就是 $X$ 的上同调环。这个答案在量子力学上会被瞬子修正所修正，修正量由我们刚才描述的那些在非常值全纯映射模空间上的积分给出。因此，手征环是对 $X$ 通常上同调环的量子修正；正因如此，手征环通常被称为 $X$ 的量子上同调环。

% [Tutor Note: 「常值映射恰好填满某一分支中的全部全纯映射」一句容易读成「常值映射构成全体全纯映射」，关键限定语是「某一分支」：$\operatorname{Maps}(\Sigma,X)$ 按同伦类分成分支，常值映射恰好填满零伦分支中的全部全纯映射。理由是全纯映射的能量 $\int_\Sigma|d\Phi|^2$ 只依赖 $\Phi_*[\Sigma]$ 所代表的同调类，在零伦类上必须为零，于是 $d\Phi\equiv0$。再与 p.220 的维数条件对照：那里 $a+b+c=n-1+dn$ 而 $\deg\sigma=2$，真正饱和虚拟维数的是 $\sum_i\deg\lambda_i=2(a+b+c)$，即带三个标记点的映射模空间的虚拟维数 $2(n-1+dn)$；文中「$=\dim\mathcal{M}_d$」若按不带标记点的模空间理解则为 $nd+n-4$，相差的 $3$ 正来自三个标记点（$\mathrm{Aut}(\mathbb{P}^1)$ 的三维）。这几个 $2$ 与 $3$ 是抄录维数条件时最常见的错误来源。]

\subsection{环结构的更多细节}
% 原文 §15.6：More Details on the Ring Structure

在给出量子上同调作为手征环的一般定义、并说明了计算它所需的数学之后，现在该更具体一些、算一个例子了。我们取 $\Sigma=S^2$ 映入紧 Kähler 流形 $X$ 的情形。先从一个只有两个算符的关联函数开始：
\[
  \langle \mathcal{O}_{\lambda_1}\mathcal{O}_{\lambda_2}\rangle,
\]
其中 $\lambda_1$ 的次数加上 $\lambda_2$ 的次数等于 $X$ 的维数。我们当然知道，常值映射的空间对这个关联函数的贡献等于
\[
  \int_X\lambda_1\wedge\lambda_2.
\]
在这种情形没有量子修正，这也就是手征环中的答案。没有量子修正的理由是：映入 $X$ 的非常值全纯映射 $S^2\to X$（并指定点 $x_1,x_2\in S^2$ 分别映到 $H_1$ 与 $H_2$）构成的空间上有一个自由的 $\mathbb{C}^*$ 作用，它由 $S^2$ 上固定 $x_1,x_2$ 的 $\mathbb{C}^*$ 作用诱导而来。该空间的某个分支要对关联函数给出非零修正，其形式维数（formal dimension）必须为零。若其实际维数为零、其上又有自由 $\mathbb{C}^*$ 作用，则它必为空集。（即便该空间可能有非空的分支，自由 $\mathbb{C}^*$ 作用也保证其贡献为零。）于是，在手征环与 $X$ 上同调的等同之下，两点关联函数——它计算手征环底下向量空间上的内积——给出上同调上通常的内积，等价地给出同调上的相交形式（intersection form）。

% ===== 原文 p.220 =====

现在来考虑环结构。取 $H^*(X)$ 的一组基 $\{\lambda_i\}$，设在这组基下内积由 $\eta_{ij}$ 给出：
\[
  \eta_{ij}=\langle \mathcal{O}_{\lambda_i}\mathcal{O}_{\lambda_j}\rangle.
\]
于是环结构由结构常数（structure constants）$c_{ij}^k$ 决定：
\[
  \mathcal{O}_{\lambda_i}\mathcal{O}_{\lambda_j}=c_{ij}^k\,\mathcal{O}_{\lambda_k}.
\]
令 $\omega_{ijk}$ 为三点函数
\[
  \omega_{ijk}=\langle \mathcal{O}_{\lambda_i}\mathcal{O}_{\lambda_j}\mathcal{O}_{\lambda_k}\rangle.
\]
由于这个三点函数也是前两个局域算符之积与第三个的内积，我们看到
\[
  \omega_{ijk}=\sum_{r}c_{ij}^r\,\eta_{rk}.
\]
由于 $\eta_{ij}$ 非退化，这说明两点与三点关联函数决定了手征环的环结构。

\subsection{$\mathbb{CP}^{n-1}$ 的计算}
% 原文 §15.7：Calculations for $\mathbf{CP}^{n-1}$

现在来看一个例子。我们借助规范线性 $\sigma$ 模型（gauged linear sigma model）来计算 $\mathbb{CP}^{n-1}$：规范群为 $U(1)$，线性玻色场为 $\phi_1,\dots,\phi_n$，如第 13 讲 §13.4 那样（那一讲中这些场记作 $A_i$）。要得到整个 $\mathbb{CP}^{n-1}$，我们令超势 $W$ 等于零。这样得到一族依赖于复参数 $t=-ir+\frac{\theta}{2\pi}$ 的理论；如我们所见，当 $\operatorname{Im}(t)\ll0$ 时，它们是 $X$ 上的非线性 $\sigma$ 模型，其 Kähler 度量依赖于 $t$，而 Kähler 类大致正比于 $\operatorname{Im}(t)$。唯一的可观测量是 $\mathcal{O}_H$，其中 $H$ 是超平面截影（hyperplane section）。我们把这个算符记作 $\sigma$。经典地，在 $\mathbb{CP}^{n-1}$ 中我们有一个关系，即 $\sigma^n=0$。让我们看看量子力学上会发生什么。我们需要计算三点函数
\[
  \langle\sigma^a\sigma^b\sigma^c\rangle,
\]
其中 $a+b+c=n-1+dn=\dim\mathcal{M}_d$，即 $d$ 次有理曲线模空间的维数。显然，只需考虑 $a,b,c<n$ 的情形。这意味着 $d$ 只可能取 $d=0,1$。当然，$d=0$ 是常值映射的模空间，其贡献就是经典的环结构。对 $d=1$，我们考虑的是 $\mathbb{CP}^{n-1}$ 中的直线。为计算
\[
  \langle\sigma^{n-1}(x)\sigma^{n-1}(y)\sigma(z)\rangle,
\]
我们必须考虑过两点的直线中与某个固定超平面相交的那些；这样的直线恰好有一条。由于瞬子作用量为 $e^{-2\pi it}$，这给出关系
\[
  \langle\sigma^{n-1}(x)\sigma^{n-1}(y)\sigma(z)\rangle=e^{-2\pi it}.
\]
从这个三点函数反解出环结构，得到量子上同调中的
\[
  \sigma^n=e^{-2\pi it}.
\]
我们也可以用物理方式做同样的计算。我们考虑的是 $\sigma$ 模型 $\mathbb{CP}^1\to\mathbb{CP}^{n-1}$。在第 13 讲 §13.6.2 中我们发现，它在红外流向 $n$ 个有能隙的真空（因为这是 $c_1(X)>0$ 的情形，

% [Tutor Note: 注意 $d=0$ 与 $d=1$ 两项的分工：$d=0$ 给经典关系 $\sigma^n=0$，$d=1$ 给关系 $\sigma^n=e^{-2\pi it}$。而「$a,b,c<n$ 已经足够」的理由是：任取 $a\ge n$，都可用 $\sigma^n=q$ 反复约化到更低次，故三点函数只依赖 $\sigma$ 的幂次模去关系——这正是「从三点函数反解环结构」的实际操作。易错点有二：(i)「过两点的直线与固定超平面相交者恰有一条」在 $n\ge3$ 时成立，但 $n=2$（即 $\mathbb{CP}^1$）时「超平面」退化为点，几何退化，故 $n=2$ 时 $\sigma^n=e^{-2\pi it}$ 的 $d=1$ 部分须另行核验；(ii) 指数上是 $e^{-2\pi it}$，符号来自 $t=-ir+\theta/2\pi$ 与 $\theta$ 项 $-i\theta\int F/2\pi$ 的约定；若某文献取 $t=\theta/2\pi+ir$，同一关系会被写成 $q=e^{2\pi it}$，对读时极易差一个整体符号。]

% ===== 原文 p.221 =====

流是从 $\sigma$ 模型流向 Landau--Ginzburg 模型）。场 $\sigma$ 在每个真空中取不同的期望值。由于诸场有质量，我们干脆把它们取为各自的期望值。这些真空期望值为 $\sigma=\mu e^{2\pi i(-t+k)/n}$，$k=0,\dots,n-1$。它们给出量子上同调环的幂等元（idempotents），由此我们推得与前面相同的答案，\footnote{除 $\mu$ 的一个幂次之外——这源于我们的物理约定与数学约定之间的不匹配。在第 13 讲中，$\sigma$ 被定义为一个量纲为一的场，且并未试图把它与任何拓扑归一化相比较。而在今天的数学讨论中，$\sigma$ 是按拓扑方式归一化的，即作为与射影空间的超平面截影相联系的场，因此其量纲为零。要比较这两个定义，应在第 13 讲的框架下，在 $r\gg0$ 的区域把 $\sigma$ 场积掉，以换得 $\sigma$ 模型的无质量场。在此过程中，$\sigma$ 将变成低能理论的一个算符，该算符可由 $\mathbb{CP}^{n-1}$ 上的一个 2-形式导出。这个算符的上同调类将是超平面截影的某一倍数 $C$，而 $C$ 正是第 13 讲的物理归一化与「给出 $\sigma^n=e^{-2\pi it}$ 的拓扑归一化」不相符的因子。（原文 p.221 页脚注）} 即该环由 $\sigma$ 生成，并服从关系 $\sigma^n=e^{-2\pi it}$。

\subsection{Fano 超曲面的计算}
% 原文 §15.8：Calculations for Fano Hypersurfaces

现在把这些计算推广到 $\mathbb{CP}^{n-1}$ 中次数 $d<n$ 的超曲面，即所谓的 Fano 超曲面。

据 [A. Collino and M. Jinzenji, \emph{On the structure of small quantum cohomology rings for projective hypersurfaces}, preprint hep-th/9611053.]，用数学方法算出的答案是：超平面截影 $\sigma$ 服从如下关系：
\begin{equation}
  \tag{15.1}
  \label{eq:15-1}
  \sigma^{d-1}\bigl(\sigma^{n-d}-e^{-2\pi it}\bigr)=0.
\end{equation}

让我们想一想，这个计算如何与下述事实相容：这个 $\sigma$ 模型在红外同样流向一个 Landau--Ginzburg 模型，再加上额外的一些有质量真空。方程 $\sigma^{d-1}\bigl(\sigma^{n-d}-e^{-2\pi it}\bigr)=0$ 的根有两类：$\sigma=0$ 处一个重数为 $d-1$ 的根，以及 $\sigma\ne0$ 处的非退化根（即重数为一的根）。方程的根对应于量子场论的真空。非退化根对应于有质量真空，它们给出量子上同调环的幂等元。在第 13 讲 §13.6.2 中分析该理论的真空结构时，我们找到 $\sigma\ne0$ 处的 $n-d$ 个有质量真空。它们位于 $e^{-2\pi it}$ 的 $(n-d)$ 次根处（再乘以一个常数，其解释与上一条脚注相同），并给出量子上同调方程中的因子 $\bigl(\sigma^{n-d}-e^{-2\pi it}\bigr)$。另一个因子 $\sigma^{d-1}$ 是原点处 Landau--Ginzburg 真空的贡献。这个因子在 $d>1$ 时出现（对 $d=1$，我们讨论的是余维为一的射影空间，没有 Landau--Ginzburg 真空）；当 $d>2$ 时，它对应于原点处一个重数大于一的根。这一限制的理由从物理上看很清楚：对 $d=2$，Landau--Ginzburg 理论的中心荷（一般情形为 $\hat c=n(1-2/d)$）为零，因此这时 Landau--Ginzburg 理论有质量、在红外是平庸的。对 $d>2$，Landau--Ginzburg 理论有一个非平凡临界点，它自然可以贡献一个重数大于一的根。由于标度不变性，这个根将位于原点。如何计算 Landau--Ginzburg 轨形（orbifold）的上同调，在其中一个习题中有描述；有了那些结果，就可以算出眼前这个情形中根的重数。

% [Tutor Note: 检验 (15.1) 是否漏项的最快办法是数根（计重数）：$(n-d)+(d-1)=n-1$，恰等于 $X$ 的上同调总维数——$X$ 是 $\mathbb{CP}^{n-1}$ 中次数 $d$ 的超曲面，复维 $n-2$，故 $H^*(X)$ 由 $1,H,\dots,H^{n-2}$ 张成，共 $n-1$ 个类，与量子化后真空（幂等元）数应当一致。若把 $\sigma^{d-1}$ 误抄成 $\sigma^{d-2}$，或把 $e^{-2\pi it}$ 的根数 $(n-d)$ 与 $n$ 搞混，这个计数立刻对不上。另可注意：原文说「$d>2$ 时该根重数大于一」，而 $d=2$ 时 $\hat c=0$、LG 理论有质量，故 $\sigma^{d-1}=\sigma$ 的根重数为一，两者自洽。]

% ===== 原文 p.222 =====

\subsection*{参考文献与交叉引用}
% 原文：References and cross references

本讲大部分内容的一个参考文献是

E. Witten, \emph{On the Landau--Ginzburg description of $\mathcal{N}=2$ minimal models}, Int.\ J.\ Mod.\ Phys.\ A \textbf{9} (1994), pp.\ 4783--4800.

关于 Landau--Ginzburg 及相关模型的椭圆亏格的进一步计算，参见

P. Berglund and M. Henningson, \emph{On the Elliptic Genus and Mirror Symmetry}, in \emph{Essays on
Mirror Manifolds II}, ed.\ S.-T. Yau (International Press), and references therein.
